%----------HEADING----------%
\begin{center}
    \textbf{\Huge \scshape Roben Bhatti} \\ \vspace{1.5pt}
    \href{mailto:robenbhatti@gmail.com}{\seticon{faEnvelope} robenbhatti@gmail.com} \quad
    \href{https://www.linkedin.com/in/roben-bhatti}{\seticon{faLinkedin} linkedin.com/roben-bhatti} \quad
    \href{https://github.com/bhroben}{\seticon{faGithub} github.com/bhroben} \quad
    \\ \vspace{2pt}
    \seticon{faPhone} \ \small +39 331 211 6804 \quad 
    %\seticon{faMapMarker} \ \small Brescia, Italy \quad 
    \seticon{faFlag} \small Italian citizenship \quad 
\end{center}


%-----------EXPERIENCE-----------%
\section{Experience}
\resumeSubHeadingListStart

    \resumeSubheading
    {German Aerospace Center (DLR)}{Oct 2024 - Current}
    {Data Scientist Intern}{Bremen, Germany}
    \resumeItemListStart
        \resumeItem{Developed and implemented Machine Learning models to quantify uncertainty in aerodynamic coefficients for reusable launch vehicles, leading to more robust estimation and performance analysis.}
        \resumeItem{Built a \textbf{CI/CD} pipeline using \textbf{Gitlab} and \textbf{Docker}, automating the deployment process and ensuring reliable releases.}
        \resumeItem{Collaborated within an \textbf{Agile team} to deliver data-driven results, culminating to the successful completion of my Master Thesis and a \textbf{co-authored} conference paper (IAC 2025) detailing the novel application (\textit{in preparation}).}
        \resumeItem{Addressed computational bottlenecks in model training by implementing Sparse Gaussian Processes, resulting in a \textbf{30\% decrease in model training time} and enhancing the feasibility of large-scale uncertainty analysis.}
    \resumeItemListEnd


\resumeSubHeadingListEnd

%-----------EDUCATION-----------%
\section{Education}
    \resumeSubHeadingListStart

    \resumeSubheading
    {University of Padua}{Oct 2022 - Jul 2025}
    {Master of Science in Physics "Physics of Data" (GPA: 4.00 / 4.00)}{Padua, Italy}
    \resumeItemListStart
        \resumeItem{\textbf{Relevant Coursework:} Mathematical and Numerical Methods, Deep Learning and Neural Networks, Advanced Statistics for Physics Analysis, Information Theory and Inference, General Relativity.}
    \resumeItemListEnd
    
    
    \resumeSubheading
    {University of Padua}{Oct 2019 - Oct 2022}
    {Bachelor of Science in Astronomy (GPA: 3.45 / 4.00)}{Padua, Italy}
    \resumeItemListStart
        \resumeItem{\textbf{Relevant Coursework:} Advanced Calculus, Statistics, Analytical Mechanics, Quantum Physics, Special Relativity.}
    \resumeItemListEnd

    \resumeSubHeadingListEnd

%-----------PROJECTS-----------%
\section{Projects}
\resumeSubHeadingListStart

    \resumeProjectHeading
    {\textbf{LLM-Powered RAG Q\&A System} (Ongoing)}{}
    \resumeItemListStart
        \resumeItem{Designed and implemented a \textbf{Retrieval-Augmented Generation (RAG)} pipeline using \textbf{LangChain} and \textbf{LangGraph}, enabling question-answering over uploaded documents (PDF, DOCX, TXT) and web content.}
        \resumeItem{Integrated Ollama for running local LLMs with \textbf{OpenAI API} fallback, returning answers with quoted citations from the original sources.}
        \resumeItem{Developed a \textbf{REST API} backend with \textbf{FastAPI}, following OOP design principles, and auto-documented endpoints via OpenAPI.}
    \resumeItemListEnd

    \resumeProjectHeading
    {\textbf{Streaming Particle Physics Data and Online Data Analysis with PySpark and Apache Kafka}}{}
    \resumeItemListStart
        \resumeItem {Engineered a distributed streaming pipeline capable of processing over \textbf{10 GB} of data per day, simulating a real-time detector stream with \textbf{Kafka} and \textbf{Spark} for live analysis.}
        \resumeItem{Deployed the environment on a cloud cluster using \textbf{AWS S3} for storage, Kafka for message queuing, and Spark for distributed processing.}
        \resumeItem{Developed a real-time interactive dashboard with Bokeh for \textbf{Data Visualization}.}
    \resumeItemListEnd

    \resumeProjectHeading
    {\textbf{Chess Position Recognition using Transformers (DETR) in PyTorch}}{}
    \resumeItemListStart
        \resumeItem{Developed a computer vision pipeline to digitize physical chess games, achieving \textbf{87.5\%} accuracy in identifying pieces and their positions from a single image.}
        \resumeItem{Fine-tuned a DETR (DEtection TRansformer) model, leveraging \textbf{transfer learning} to significantly reduce training time and improve performance on custom-built dataset.}
        \resumeItem{Developed a robust \textbf{post-processing} module to automatically convert the model's raw detection output into standard notation (FEN), making the board state immediately compatible with chess engines and analysis tools.}
    \resumeItemListEnd


    

\resumeSubHeadingListEnd

%-----------PROGRAMMING SKILLS-----------%

\section{Technical Skills}
    \begin{itemize}[leftmargin=0.15in, label={}]
	\small{\item{
		\textbf{Languages}{: Python, R, SQL} \\
		\textbf{Technologies}{: Docker, Git, CI/CD, Kafka, Spark, Pytorch, SciPy, NumPy, Pandas, LangChain, FastAPI, Tableau.} \\}}
    \end{itemize}

\section{Languages}
     \begin{itemize}[leftmargin=0.15in, label={}]
 	\small{\item{
 		\textbf{English}: Professional, \textbf{Italian}: Native }}
     \end{itemize}

\clearpage


